\documentclass[11pt]{article}
\usepackage{amsmath,amssymb,amsthm,hyperref}
\newtheorem{proposition}{Proposition}
\title{Full-Cube Counting Audit and an Unconditional Baseline}
\author{}
\date{November 23, 2026}
\begin{document}
\maketitle

\paragraph{Status.}
Neither the requested upper bound $O(2^n/(n\log\log n))$ nor the
alternative lower bound $\Omega(2^n/n)$ is proved by this note.
The unconditional result imported and checked here is the existing
$\Omega(2^n/n^2)$ counting baseline. The subsequent calculations record
what a sharper full-cube argument would have to improve. They are not
asserted as an obstruction to all admissible representations.

\section{Checked unconditional counting theorem}
The public module
\texttt{HeadComplexity/Results/WarrenBound.lean} and its twelve new
polynomial dependencies are copied verbatim from
\texttt{VerifiedMechanisms/xorformer-runs}, commit
\texttt{6f5c7d43e32195fe4b05508de9f5f2d3e4888699}.
The imported dependency closure has 1,450 lines. It proves Warren's
strict sign-pattern estimate from algebraic root counts, generic regular
values, and constrained minima. No sign-pattern theorem is assumed.
The current files defining the model, \texttt{HStar}, and
\texttt{worstCaseHeadComplexity} are unchanged.

The declarations
\texttt{warrenStrictSignPatternBound},
\texttt{card\_headComputableFunctions\_le\_counting}, and
\texttt{worstCaseHeadComplexity\_gt\_div}
give, respectively, the proved Warren interface, the bound
\[
 \#\{f:H^*(f)\leq H\}\leq2^{35n^2H}
 \quad(1\leq n,\ 1\leq H\leq2^n),
\]
and the unconditional finite statement
\[
 \left\lfloor\frac{2^n-1}{35n^2}\right\rfloor<W(n)
 \quad\text{when }1\leq n\text{ and }35n^2\leq2^n-1.
\]
The floor in the last display is natural-number division in the Lean
declaration. It yields the stated asymptotic baseline.
This does not supply the extra factor of $n$ sought in the
alternative lower-bound target.

\paragraph{Verification.}
\texttt{lake build HeadComplexity.Results.WarrenBound} passed with
3,426 jobs and no warnings. Independent \texttt{\#print axioms} checks of
all six public declarations in the entrypoint reported only
\texttt{propext}, \texttt{Classical.choice}, and \texttt{Quot.sound}.
The imported files contain no \texttt{sorry} and no new axiom.
The build and axiom output is preserved in
\texttt{reverse-alternative-warren-check.txt}.

\section{Exact full-cube coefficient recurrence}
For $x\in\{0,1\}^n$, let
$A(x)=a_0+\sum_i a_ix_i$ and $B(x)=b_0+\sum_i b_ix_i>0$.
Write the unique multilinear expansion of $A/B$ as
$\sum_{S\subseteq[n]}u_S x_S$, where $x_S=\prod_{i\in S}x_i$.
\begin{proposition}
For every $S\subseteq[n]$,
\[
 \left(b_0+\sum_{i\in S}b_i\right)u_S
       +\sum_{i\in S}b_i u_{S\setminus\{i\}}=a_S,
 \qquad
 a_S=\begin{cases}a_0&S=\varnothing,\\a_i&S=\{i\},\\0&|S|\geq2.
 \end{cases}
\]
\end{proposition}
\begin{proof}
Multiply the multilinear expansion by $B$ in the quotient algebra
$\mathbb R[x_1,\ldots,x_n]/(x_i^2-x_i)_{i=1}^n$.
The coefficient of $x_S$ in $b_ix_i\sum_Tu_Tx_T$ is
$b_i(u_S+u_{S\setminus\{i\}})$ if $i\in S$, and zero otherwise.
The resulting function is $A$, whose coefficients are the displayed $a_S$.
\end{proof}
The diagonal coefficient is $B(\mathbf1_S)>0$, so this is a genuine
triangular recurrence even with the negative denominator orientation.
For a sum of heads the recurrences have different coefficients.
Composing their multiplication operators returns the familiar
degree-$H$ cleared numerator. No stronger recurrence for their sum, or
sign-only lower bound, has been obtained from this calculation.

\section{Convex sections and the remaining multiplicity issue}
Fix all heads except one and fix a required sign function $f$ on the cube.
Let $R(x)$ be the score of the fixed heads and the bias. On the region
$B(x)>0$, the sign constraints on the remaining head are exactly
\[
 f(x)\bigl(A(x)+R(x)B(x)\bigr)>0\quad\text{for every }x.
\]
These are linear inequalities in that head's affine coefficients.
For either fixed denominator orientation, admissibility is also a
collection of linear strict inequalities. Thus every such one-head
section of the feasible parameter set is convex.

This observation does not imply that the complete feasible set is
connected or that its intersections have few isolated solutions.
The checked Warren proof bounds the relevant constrained minima by
$(2d)^p$. With $p=1+2H(n+1)$ and $d=H+1$, this is still
$H^{O(Hn)}$. Replacing it by a suitable full-cube bound on the scale
$H!C^{Hn}$ would be quantitatively useful when $H$ is exponential in
$n$, but no such result is proved here. A uniform estimate of that
strength for arbitrary promised subsets would conflict with the
existing positive-orientation one-hot decoder. Any argument must
specify and use its full-cube or parameter-regime hypotheses.

\section{A sparse-polynomial shortcut does not follow}
There is a simple exact interpolation construction that cautions against
transferring Fourier-sparsity lower bounds directly to heads.
Partition the input into two blocks $x,y\in\{0,1\}^k$ and set
$t(x,y)=|x|+(k+1)|y|$. It takes the $(k+1)^2$ distinct values
$0,\ldots,(k+1)^2-1$.
For $K=(k+1)^2$ choose distinct positive numbers $p_1,\ldots,p_K$.
The Cauchy matrix $[1/(t+p_h)]_{0\leq t<K,1\leq h\leq K}$ is invertible.
Indeed a linear dependence, multiplied by $\prod_h(t+p_h)$,
would give a polynomial of degree at most $K-1$ vanishing at $K$
distinct points. The polynomial is zero; evaluating at $-p_h$ then
shows that every coefficient of the dependence is zero.

Consequently every real function of the two block weights has an exact
$K$-head representation with constant numerators and denominators
$p_h+t(x,y)$. These denominators have strictly positive slopes and
positive constants, so they are actual original-model atoms.
In particular the intersection of two majority functions has a
polynomial head bound by this elementary construction. Sherstov's
Theorem 6.3 gives Fourier threshold density
$k^{\Omega(\log k)}$ for that function; see
\url{https://web.cs.ucla.edu/~sherstov/pdf/hshs.pdf}.
Thus a general conversion to a Fourier polynomial with only
$H\,\operatorname{poly}(n)$ monomials cannot be used here.
This literature-based consequence is not Lean-formalized and is not
used by the imported counting theorem.

\section{Positive poles are not uniquely recoverable from the cube}
There is a nontrivial exact equality between two sums of three positive
reciprocal heads. Let $a=x_1-x_2$, $b=x_2-x_3$, and let $D$ be any
affine form with positive constant and all slopes strictly greater than
two. Then on the entire Boolean cube
\[
 \frac1{D+2a+2b}+\frac1{D+a}+\frac1{D+b}
 =\frac1{D+2a+b}+\frac1{D+a+2b}+\frac1D.
\]
Indeed at every Boolean point one of $a=0$, $b=0$, or $a+b=0$ holds,
since among three bits two are equal. In each case the two denominator
triples coincide as multisets. Every displayed denominator retains
positive constant and strictly positive slopes. The triples are
different affine-form multisets, and every numerator is the constant one.

The ambiguity is not explained merely by a continuous redundant fiber.
For dimension five, take
$D=100+4x_1+5x_2+6x_3+7x_4+8x_5$ and view each triple as a map from
its eighteen independent affine denominator coefficients to cube values.
Both Jacobians have column rank eighteen. An exact certificate uses
the eighteen lexicographic vertex indices
\[
 0,16,8,4,3,5,6,17,9,12,10,11,2,18,14,15,7,13.
\]
The two square minors have determinants congruent to $2363$ and
$997640$, respectively, modulo the prime $1000003$. Every denominator
in the Jacobian is invertible modulo that prime. Thus both rational
determinants are nonzero. The reproducible exact calculation is
\texttt{reverse-alternative-pole-check.py}. The inverse function theorem
applied to those selected coordinates shows that each representation is
locally isolated in its full-cube value fiber. This regularity consequence
is written, not Lean-formalized.

The example refutes uniqueness of positive-pole representations even
with constant positive numerators and local isolation. It yields only
two distinct unordered representations. It does not refute the much
weaker quantitative bound $H!C^{Hn}$ discussed above and does not give
an asymptotic lower bound or an improved compiler.

\paragraph{Open work.}
The exact recurrence, convex sections, shared-direction observations,
and standard rational-network counting results have not supplied the
needed positivity-sensitive multiplicity estimate. No universal
global construction achieving the logarithmic saving has been found
by this route either.
\end{document}
